% !TEX root = ../Main.tex
\chapter{压强与力}

压强与力的关系是 $F=pS$：一块面积 $S$ 上若受到均匀压强 $p$，它受到的压力就是 $F=pS$。下面用这条关系，从一个"看似要算、其实抵消"的例子讲起。

\section{大气压：上下抵消}

人站在地面上，头顶受大气向下压、脚底受地面向上支持。先问一个常见的问题：脚底受到的支持力 $N$ 等于体重 $G$ 吗？

\begin{center}
\begin{tikzpicture}[>=Stealth, scale=1]
    \draw (-1.6,0) -- (1.6,0); % 地面
    \draw (0,1.0) circle (0.32); % 头
    \fill (-0.1,1.08) circle (0.5pt); \fill (0.1,1.08) circle (0.5pt);
    \draw (0,0.68) -- (0,-0.2); % 身体
    \draw (-0.5,0.45) -- (0.5,0.45); % 手臂
    \draw (0,-0.2) -- (-0.35,-0.7); \draw (0,-0.2) -- (0.35,-0.7); % 腿
    \draw[->, thick] (0,1.32) -- (0,1.95) node[right] {$N$};
    \draw[->, thick] (0,0.4) -- (0,-1.0) node[right] {$G$};
\end{tikzpicture}
\end{center}

如果只盯着体重，会以为 $N=G$。但别忘了大气压。从上往下看，人占的那块地面面积上方，整根空气柱压下来；这块面积下方，地面又顶上来。把大气压也算进去：

\[
N=G+p_0 S_{\text{上}}\quad?\qquad\text{——还要看大气压在下表面有没有抵消。}
\]

\begin{center}
\begin{tikzpicture}[>=Stealth, scale=1]
    \draw (-2.4,-1.5) rectangle (2.4,1.5);
    \node at (-2.4,1.75) {从上往下看};
    \draw (0,0) circle (1.1);
    \draw[->, very thick] (0.25,0.55) -- (-0.15,-0.05) node[midway, right] {$pS$};
    \draw[->, thick] (-0.65,-0.75) -- (-0.95,-1.05) node[below right] {$pS$};
\end{tikzpicture}
\end{center}

关键在于：从整体上看，人占据的截面，\textbf{上表面}面积与\textbf{下表面}面积是相等的，

\[
S_{\text{上}}=S_{\text{下}}\quad\Longrightarrow\quad p_0 S_{\text{上}}=p_0 S_{\text{下}}.
\]

大气从上方压下来的力 $p_0 S_{\text{上}}$，正好被从下方（地面经接触面）顶上来的力 $p_0 S_{\text{下}}$ 抵消。\textbf{总效果为"无"}——大气压对竖直方向的合力没有贡献，所以最后仍是 $N=G$。

\rmk{这就是为什么平时算支持力可以不管大气压：上下表面压力相等、自动抵消。}

\section{又如：浸在液体中的物体}

同样的"上下表面压力差"思路，用到液体里就是浮力。

\ex{把一个正方体悬浮在水中，画出它的受力示意图。}

\sol{
    正方体上表面受水向下的压力 $F_{\text{上}}$，下表面受水向上的压力 $F_{\text{下}}$，自身重力 $G$ 向下。

    \begin{center}
    \begin{tikzpicture}[>=Stealth, scale=1]
        \draw (-2.4,-1.7) rectangle (2.4,1.9); % 容器
        \foreach \y in {-1.3,-0.8,-0.3,0.2,0.7,1.2} \draw[dashed, gray] (-2.4,\y) -- (2.4,\y); % 水位线
        \draw (-0.6,-0.6) rectangle (0.6,0.6); % 正方体
        \draw[->, thick] (-0.3,0.6) -- (-0.3,1.05) node[above] {$F_{\text{上}}$};
        \draw[->, thick] (0.3,-0.6) -- (0.3,-1.05) node[below] {$F_{\text{下}}$};
        \draw[->, thick] (0,-0.6) -- (0,-1.4) node[below] {$G$};
    \end{tikzpicture}
    \end{center}

    下表面比上表面深，水压更大，所以 $F_{\text{下}}>F_{\text{上}}$，二者之差就是浮力：
    \[
    F_{\text{浮}}=F_{\text{下}}-F_{\text{上}}.
    \]
}

\ex{把一个正方体沉在容器底部（与底紧贴、底面无水），画出它的受力并求底部对它的支持力 $N$。}

\sol{
    沉底后底面紧贴容器，水进不到下表面，于是下方不再有水的向上压力，取而代之的是容器底的支持力 $N$。上表面仍受水向下的压力 $F_{\text{上}}=p_0 s$（这里 $p_0$ 是上表面处水压，$s$ 为底面积），重力 $G$ 向下。

    \begin{center}
    \begin{tikzpicture}[>=Stealth, scale=1]
        \draw (-2.4,-1.7) rectangle (2.4,1.9);
        \foreach \y in {-1.3,-0.8,-0.3,0.2,0.7,1.2} \draw[dashed, gray] (-2.4,\y) -- (2.4,\y);
        \draw (-0.6,-1.7) rectangle (0.6,-0.5); % 贴底正方体
        \draw[->, thick] (-0.3,-0.5) -- (-0.3,0.0) node[above] {$F_{\text{上}}$};
        \draw[->, thick] (0,-0.5) -- (0,-0.05); \node[right] at (0,-0.25) {$G$};
        \draw[->, thick] (0.3,-1.7) -- (0.3,-2.15) node[below] {$N$};
        \draw[->] (2.7,-1.7) -- (2.7,1.9) node[midway, right] {$h$};
    \end{tikzpicture}
    \end{center}

    竖直方向受力平衡，容器底给的支持力要顶住重力和上方水压：
    \[
    N=F_{\text{上}}+G=mg+\rho g h s.
    \]
}

\rmk{沉底"贴死"是关键：一旦下表面没有水，浮力的"下表面向上压力"这一项就消失了，受力图和悬浮时完全不同。}

\section{思考}

\textbf{1. 浮力是怎么产生的？}\quad 就是上下表面的压力差：
\[
F_{\text{浮}}=F_{\text{下}}-F_{\text{上}}.
\]
把上下表面的水压逐层算清楚，就能推出\textbf{阿基米德原理}
\[
F_{\text{浮}}=G_{\text{排}},
\]
即浮力等于被排开的那部分液体的重力。

\textbf{2. 把一个半球贴底放在容器底部（平面圆底朝下贴在容器底、球面朝上凸入水中，形如锅盖），它受到容器底部的支持力有多大？}

\sol{
    核心是\textbf{假设法}：假设半球\textbf{下方也有水}（半球完全浸没）。这样半球就像普通浸没物体一样受到完整的浮力，
    \[
    F_{\text{浮}}'=\rho V g=\rho\cdot\tfrac{2}{3}\pi R^3\cdot g,
    \]
    其中 $V=\tfrac{2}{3}\pi R^3$ 是半球体积。完整浮力 $=$ 平底（深度 $h$ 的圆面）受到的向上水压 $-$ 球面受到的向下水压，而平底向上水压为 $\rho g h\,\pi R^2$，于是球面受到的向下水压为
    \[
    F_{\text{下}}=\rho g h\,\pi R^2-F_{\text{浮}}'=\rho g h\,\pi R^2-\tfrac{2}{3}\pi R^3\rho g.
    \]
    回到真实情形：半球贴底，下方其实没有水，平底那份向上水压不存在，改由容器底的支持力顶住；球面那份向下水压 $F_{\text{下}}$ 依旧在。容器底要顶住的正是这份向下水压，
    \[
    N=F_{\text{下}}=\rho g h\,\pi R^2-\tfrac{2}{3}\pi R^3\rho g=\pi\rho g R^2\Big(h-\tfrac{2}{3}R\Big).
    \]
}

\begin{center}
\begin{tikzpicture}[>=Stealth, scale=1]
    \draw (-2.4,-1.2) rectangle (2.4,1.6);
    \foreach \y in {-0.8,-0.3,0.2,0.7,1.2} \draw[dashed, gray] (-2.4,\y) -- (2.4,\y);
    \fill[gray!12] (0.9,-1.2) arc (0:180:0.9) -- cycle;
    \draw (0.9,-1.2) arc (0:180:0.9); % 球面朝上（穹顶）
    \draw (-0.9,-1.2) -- (0.9,-1.2); % 平面朝下贴底
    \node at (1.2,-1.0) {$R$};
    \draw[->, thick] (0,0.0) -- (0,-0.55) node[below right] {$F_{\text{下}}$};
    \draw[->] (2.7,-1.2) -- (2.7,1.6) node[midway, right] {$h$};
\end{tikzpicture}
\end{center}

\rmk{"假设法"的妙处：先补上不存在的水，把问题凑成会算的浮力，再把多补的那份减掉。}
